% Part: incompleteness
% Chapter: incompleteness-provability
% Section: provability-conditions

\documentclass[../../../include/open-logic-section]{subfiles}

\begin{document}

\olfileid{inc}{inp}{prc}

\olsection{د~$\Th{PA}$ د \usetoken{S}{derivability} شرطونه}

د پېانو حساب، يا~$\Th{PA}$، هغه تيوري ده چې~$\Th{Q}$ د ټولو
!!{formula}s دپاره د استقرا د بديهي اصولو په زياتولو غځوي۔ په نورو
ټکو، $\Th{Q}$ ته د دې بڼې بديهي اصول ورزياتېږي:
\[
(!A(0) \land \lforall[x][(!A(x) \lif !A(x'))]) \lif \lforall[x][!A(x)]
\]
د هر !!{formula}~$!A$ دپاره۔ پام وکړئ چې دا په حقيقت کښې يوه
\emph{شېما} ده، يعنې بې‌شمېره بديهي اصول دي (او دا هم څرګندېږي
چې~$\Th{PA}$ په متناهي ډول {\em نۀ} شي بديهي کېدے)۔ خو ځکه چې
په مؤثره طريقه معلومولے شو چې د نښو يوه لړۍ د استقرا د کوم بديهي
اصل بېلګه ده که نۀ، د~$\Th{PA}$ د بديهي اصولو سټ محاسبه کېدونکے
دے۔ $\Th{PA}$ له~$\Th{Q}$ څخه ډېره پياوړې تيوري ده۔ مثلاً د
عادي استقرا په کارولو په آسانۍ ثابتولے شو چې جمع او ضرب تبادلي
دي۔ په حقيقت کښې د عددونو د تيورۍ او ترکيبپوهنې زياتره
متناهي‌پاله استدلالونه په~$\Th{PA}$ کښې ترسره کېدے شي۔

څرنګه چې~$\Th{PA}$ په محاسبوي ډول بديهي‌کړې ده، د
!!{derivability} پريديکات~$\Prf[\Th{PA}](x,y)$ محاسبه کېدونکے
دے او ځکه په~$\Th{Q}$ کښې، او په دې توګه په~$\Th{PA}$ کښې هم،
تمثيلېږي۔ د پخوا په شان د دې اړيکې د تمثيلوونکي فورمول دپاره
$\OPrf[\Th{PA}](x,y)$ کاروو۔ $\OProv[\Th{PA}](y)$ دې فورمول
$\lexists[x][\OPrf[\Th{PA}](x,y)]$ وي؛ دا په مفهومي ډول وايي:
«$y$ د~$\Th{PA}$ له بديهي اصولو څخه !!{derivable} دے»۔
\textbf{د ژباړې د نسخې سمون (OLPRV-001):} په اصلي سرچينه کښې د
همدې تعريف دننه~$\Prf[\Th{PA}](x,y)$ ليکل شوے، خو دا د عددونو
تر منځ بهرنۍ اړيکه ده؛ دلته د هغې تمثيلوونکے د ژبې فورمول
$\OPrf[\Th{PA}](x,y)$ پکار دے، لکه چې مخکينۍ جمله ئې په ښکاره
نوموي۔ له~$\Th{Q}$ د بديهي اصولو څخه لږ څه زيات ته ځکه اړ يو
چې خپله تيوري بايد دومره پياوړې وي چې د دې !!{derivability}
پريديکات په باب څو بنسټيز حقايق !!{derive} کړي۔ اړين حقايق دا دي:
\begin{enumerate}
\item[P1.] که $\Th{PA} \Proves !A$ وي، نو $\Th{PA} \Proves
  \OProv[\Th{PA}](\gn{!A})$۔
\item[P2.] د ټولو !!{formula}s $!A$ او $!B$ دپاره:
  \[
  \Th{PA} \Proves \OProv[\Th{PA}](\gn{!A \lif !B}) \lif
  (\OProv[\Th{PA}](\gn{!A}) \lif \OProv[\Th{PA}](\gn{!B})).
  \]
\item[P3.] د هر !!{formula}~$!A$ دپاره:
  \[
  \Th{PA} \Proves \OProv[\Th{PA}](\gn{!A})
  \lif \OProv[\Th{PA}](\gn{\OProv[\Th{PA}](\gn{!A})}).
  \]
\end{enumerate}
دا درې خاصيتونه يوازې په دې توګه کتلے شو چې !!{formula}
$\OProv[\Th{PA}](y)$ په دقت بيان کړو او د~$\Th{PA}$ بديهي اصول
د اړوندو صوري !!{derivation}s د بيانولو دپاره وکاروو۔ شرطونه (۱)
او~(۲) آسانه دي؛ اصلي کار شرط~(۳) غواړي۔ (فکر وکړئ چې دا به
څه ډول کار وي \dots) د تفصيلاتو بشپړول ستړي کوونکي او دلته
بې‌ګټې وي، نو اوس له تاسو غواړو چې دا ومنئ چې~$\Th{PA}$ همدا
درې پورته خاصيتونه لري۔ د~$\OProv[\Th{PA}](y)$ يوه مناسبه ټاکنه
به دا هم پوره کوي:
\begin{enumerate}
\item[P4.] که $\Th{PA} \Proves \OProv[\Th{PA}](\gn{!A})$ وي، نو
  $\Th{PA} \Proves !A$۔
\end{enumerate}
خو موږ به دې حقيقت ته اړ نۀ شو۔

\begin{digress}
په ضمن کښې ګوډل هم دلته زموږ په شان د تفصيلاتو په ليکلو کښې
تنبلي وکړه۔ د خپلې ۱۹۳۱ز مقالې په پاې کښې ئې د ناتکميلۍ د دويمې
قضيې د ثبوت لنډه نقشه ورکړه او ژمنه ئې وکړه چې تفصيلات به په بله
مقاله کښې راوړي۔ بيا ورته وخت ورنۀ غے؛ ځکه هر څوک چې په استدلال
پوهېدل، باور ئې درلود چې تفصيلات بشپړېدے شي (او هغه اړ نۀ و چې
پخپله ئې وليکي۔)
\end{digress}

\end{document}
